% !PRINTED 201
\weberpage{201}
Wir nehmen irgend eine der Gleichungen (1)

\weq{8}{a = \Theta(a)}

und setzen in der Gleichung

\weq{9}{a_1 = \psi_1(a)}

für $a$ den ihm gleichen Werth $\Theta(a)$. So erhalten wir

\weq{10}{a_1 = \psi_1 \Theta(a).}

Nach (9) ist aber $a = \psi_1^{-1}(a_1)$, und folglich ergiebt sich aus (10):

\weq{11}{a_1 = \psi_1 \Theta \psi_1^{-1}(a_1).}

Daraus ersieht man, da wir $\nu - 1$ verschiedene Substitutionen $\Theta$ haben, dass $a_1$ ein $\nu$-zähliger Pol ist und zwar zu der mit $Q$ conjugirten Gruppe $\psi_1 Q \psi_1^{-1}$ gehörig. Aus diesem Grunde nennen wir

\weq{12}{a, a_1, a_2, \ldots, a_{\mu-1}}

\gesperrt{ein System conjugirter $\nu$-zähliger Pole von $G$.}

Eine beliebige Substitution $\chi$ von $G$ kann nach (6) immer in die Form $\psi_i \Theta$ gebracht werden, so dass $\Theta$ zu $Q$ gehört, und daraus folgt, dass $\chi(a) = \psi_i \Theta(a) = a_i$ ist, und ebenso, wenn man $\psi_i \Theta \psi_k = \psi_h \Theta'$ setzt, so dass auch $\Theta'$ zu $Q$ gehört:

\[ \chi(a_k) = \psi_i \Theta \psi_k(a) = \psi_h \Theta'(a) = a_h. \]

Da ferner zwei Grössen $\chi(a_h), \chi(a_k)$ nur dann einander gleich sein können, wenn $a_h = a_k$ ist, so folgt der Satz:

5. \gesperrt{Die beiden Reihen}

\[ a, \quad a_1, \quad a_2, \quad \ldots, \quad a_{\mu-1} \]
\[ \chi(a), \quad \chi(a_1), \quad \chi(a_2), \quad \ldots, \quad \chi(a_{\mu-1}) \]

\gesperrt{stimmen, welche Substitution aus $G$ auch für $\chi$ genommen werden mag, abgesehen von der Reihenfolge, mit einander überein.}

Ist dann $b$ ein in dem Systeme (12) nicht enthaltener Pol, so kann auch $\chi(b)$ nicht in (12) vorkommen, und daraus ergiebt sich, dass zwei Systeme conjugirter Pole, wenn sie nicht ganz identisch sind, keinen Pol gemein haben.

Hiernach können wir die sämmtlichen Pole der Gruppe $G$ in Systeme conjugirter Pole anordnen, und wir erhalten nach \S.\,51 ihre Gesammtzahl $2n - 2$, wenn wir jeden $\nu$-zähligen Pol $(\nu - 1)$mal mitrechnen.
